% Source: https://artofproblemsolving.com/community/c3h2419073_telescoping_series
% Stable identity: aops-topic-2419073-q01
\textbf{(a)}
\[f(x)=\dfrac{1\cdot2\cdot3\cdots x}{2\cdot3\cdot4\cdots(x+1)}=\boxed{\dfrac{1}{x+1}},\]
since every factor $2,3,\dots,x$ appears once on top and once at the bottom.

\textbf{(b)} In every finite product the cancellation leaves two factors: the first numerator $1$ and the \emph{last} denominator $x+1$. The student cancelled all the denominators, but an infinite product has no last factor; the leftover denominator does not disappear, it only moves further out and grows. The value of an infinite product is defined as the limit of its partial products, and
\[\lim_{x\to\infty}f(x)=\lim_{x\to\infty}\dfrac{1}{x+1}=\boxed{0}.\]
This agrees with the observation that $f(x)$ keeps getting smaller. (The same faulty cancellation applied to $\dfrac21\cdot\dfrac32\cdot\dfrac43\cdots$ would also give $1$, although its partial products are $x+1$, which grow without bound.)
